\section{模型评价与推广}

\subsection{模型优点}

\begin{enumerate}
    \item \textbf{多方法交叉验证。}
    在网时长的首要地位被四种独立方法共同确认：
    问题一的 Pearson 相关系数、问题二的逻辑回归系数和 SHAP 值，
    以及第七章的特征消融实验（移除后在网时长 AUC 降幅 1.3\%）。
    同一特征在不同分析框架下均指向相同结论，降低了单方法偶然性。

    \item \textbf{以实验证据支撑线性假设。}
    本文在问题二中将逻辑回归与随机森林、XGBoost 进行了 AUC 对比
    （差值 $< 0.01$），证实该数据集上特征与流失概率之间以线性
    关系为主导。模型选型不再是"因可解释性而用线性"的先验判断，
    而是先比较、后确认的实证过程。

    \item \textbf{类别不平衡的系统性应对。}
    问题一在数据探索中即给出了流失率 26.54\% 的量化结论；
    问题二在评估阶段弃用准确率，改用 AUC 和 F1-score；
    训练阶段设置 \texttt{class\_weight='balanced'} 以损失权重
    平衡少数类的影响。召回率达到 0.78，最终能够在不误判多数
    客户的前提下提前识别近八成的潜在流失客户。
\end{enumerate}

\subsection{模型缺点}

\begin{enumerate}
    \item \textbf{线性模型无法捕捉特征间交互。}
特征消融实验发现，光纤服务的独立区分力有限（AUC 仅微增
0.1\%）。但"光纤 + 月付 + 高月费"的三元组合可能对流失
    有协同增强效应，逻辑回归的加性结构无法自动建模此类关系。
    后续可引入显式构造的交叉特征（如在网时长 $\times$ 月付合同），
    在保留线性模型可解释性的同时增强对组合风险模式的识别能力。

    \item \textbf{缺少时序行为特征。}
    当前模型基于截面数据，未纳入消费金额变动趋势、服务套餐切换
    历史等动态信息。客户流失通常是一个渐变过程而非突然事件——
    消费下降、投诉增多、使用频率降低往往先于离网——
    动态特征有助于区分"下月即将流失"和"半年后可能流失"的客户，
    为挽留动作提供更精细的时间窗口。
\end{enumerate}

\subsection{模型改进}

针对上述局限，两个方向值得进一步探索：

一是构造显式特征交互项。将 SHAP 分析确认重要且业务上自然关联
的特征对（如在网时长 $\times$ 月付合同、光纤 $\times$ 月费）
以显式交互变量的形式加入逻辑回归。这在不改变模型框架的前提下
提升了对组合风险模式的敏感度，且系数仍可逐项解释。

二是引入时序行为特征。若后续可获得历史账单数据，可构造
"近 3 月月费变化率""近 6 月服务退订次数"等动态变量，
将静态截面模型升级为时间维度的预警系统。

\subsection{模型推广}

本文的分析管线——从特征画像到预测模型、从 SHAP 归因到分位分层、
最后以 ROI 评估策略的可行性——不完全依赖于电信行业的数据结构。

在银行业信用卡流失场景中，持卡时长可类比于在网时长，
账单分期签约可类比于合同类型。互联网订阅服务（SaaS、流媒体）
同样面临月付/年付的订阅模式选择与持续付费意愿问题，
CLTV 和 ROI 的评估框架可以直接复用。两类场景在成本结构上有所
不同——银行业需计入违约风险成本，互联网服务的边际服务成本
接近零——因此 CLTV 的取值口径和经济门槛需要分别校准，
但"分层—策略—ROI"的推理链条不受影响。
